\begin{table}[htbp]
\centering
\footnotesize
\caption{Análise de Robustez e Sensibilidade da PIP — Desfecho: Casos}
\label{tab:sens_casos}
\begin{tabular}{lccccr}
\toprule
\textbf{Covariável} & \textbf{Principal} & \textbf{Ampliado} & \textbf{Em Nível} & \textbf{Sem Fortaleza} & \textbf{$N \ge 0{,}50$} \\
\midrule
Índice de Desenvolvimento Municipal & 0,997 & 0,996 & 0,912 & 0,997 & 4 \\
Leitos vinculados ao SUS & 0,989 & 0,977 & 0,910 & 0,985 & 4 \\
Município integrante de região metropolitana & 0,981 & 0,970 & 0,986 & 0,969 & 4 \\
Internações por doenças do aparelho circulatório & 0,757 & 0,739 & 0,632 & 0,724 & 4 \\
IVS Renda e Trabalho & 0,734 & 0,649 & 0,854 & 0,710 & 4 \\
Proporção da população com 60 anos ou mais & 0,699 & 0,729 & 0,738 & 0,626 & 4 \\
População beneficiária do Programa Bolsa Família & 0,636 & 0,589 & 0,362 & 0,593 & 3 \\
Internações por asma & 0,576 & 0,521 & 0,373 & 0,534 & 3 \\
Taxa de homicídios & 0,566 & 0,476 & 0,342 & 0,525 & 2 \\
Cobertura urbana de abastecimento de água & 0,525 & 0,483 & 0,315 & 0,474 & 1 \\
IDEB dos anos iniciais do ensino fundamental & 0,514 & 0,508 & 0,550 & 0,478 & 3 \\
Município integrante do semiárido & 0,507 & 0,447 & 0,593 & 0,496 & 2 \\
Unidades de saúde vinculadas ao SUS & 0,504 & 0,472 & 0,702 & 0,466 & 2 \\
IDEB dos anos finais do ensino fundamental & 0,499 & 0,451 & 0,536 & 0,465 & 1 \\
Produto Interno Bruto per capita & 0,490 & 0,445 & 0,324 & 0,471 & 0 \\
Proporção da população em extrema pobreza & 0,488 & 0,424 & 0,455 & 0,456 & 0 \\
IDEB do 3º ano do ensino médio & 0,446 & 0,403 & 0,428 & 0,407 & 0 \\
Densidade demográfica & 0,437 & 0,395 & 0,386 & 0,289 & 0 \\
Proporção da população em empregos formais & 0,431 & 0,394 & 0,312 & 0,383 & 0 \\
População estimada & 0,430 & 0,385 & 0,375 & 0,299 & 0 \\
IVS Infraestrutura Urbana & 0,426 & 0,390 & 0,338 & 0,406 & 0 \\
Internações por diabetes mellitus & 0,413 & 0,349 & 0,400 & 0,379 & 0 \\
Profissionais de saúde vinculados ao SUS & 0,410 & 0,356 & 0,392 & 0,370 & 0 \\
Proporção da população rural & 0,408 & 0,365 & 0,590 & 0,356 & 1 \\
Proporção de mulheres eleitas & 0,386 & 0,353 & 0,344 & 0,341 & 0 \\
Cobertura média de imunização em menores de um ano & 0,382 & 0,357 & 0,398 & 0,350 & 0 \\
Internações por doenças do aparelho respiratório & 0,365 & 0,328 & 0,318 & 0,323 & 0 \\
Despesa municipal com educação e cultura & 0,355 & 0,320 & 0,317 & 0,316 & 0 \\
Proporção de votos válidos do prefeito eleito & 0,346 & 0,309 & 0,309 & 0,300 & 0 \\
Índice de Desenvolvimento Humano Municipal & 0,345 & 0,306 & 0,343 & 0,302 & 0 \\
IVS Capital Humano & 0,343 & 0,304 & 0,356 & 0,302 & 0 \\
Despesa municipal com saúde e saneamento & 0,339 & 0,306 & 0,348 & 0,302 & 0 \\
Consumo de energia elétrica por 100 mil habitantes & 0,337 & 0,302 & 0,310 & 0,294 & 0 \\
Área territorial & 0,330 & 0,295 & 0,324 & 0,285 & 0 \\
Alinhamento partidário do prefeito com o governo federal & 0,329 & 0,301 & 0,305 & 0,294 & 0 \\
Auxílio emergencial & - & 0,361 & - & - & 0 \\
Transferências federais para combate à COVID-19 & - & 0,355 & - & - & 0 \\
\bottomrule
\end{tabular}
\begin{tablenotes}\footnotesize
\item \textit{Nota:} Colunas representam a Probabilidade de Inclusão Posterior (PIP) em cada especificação metodológica.
\end{tablenotes}
\end{table}